\documentclass[12pt,a4paper]{article}
\usepackage[utf8]{inputenc}
\usepackage[T1]{fontenc}
\usepackage{amsmath, amssymb}
\usepackage{listings}
\usepackage{geometry}

% Ustawienia marginesów
\geometry{margin=2.5cm}

% Ustawienia dla kodu źródłowego
\lstset{
	basicstyle=\ttfamily\small,
	breaklines=true,
	frame=single,
	captionpos=b
}

\title{Obliczanie pierwiastków zespolonych wielomianu o współczynnikach rzeczywistych
	metodą Bairstowa}
\author{Adam Kwiatkowski}
\date{Numer indeksu: 167957}

\begin{document}
	
	\maketitle
	
	\section{Zastosowanie}
	% Punkt 1: Jedno zdanie nt. zastosowania
	Procedura bairstow\_method oblicza pierwiastki zespolone wielomianu (1) o współczynnikach rzeczywistych w zmiennopozycyjnej arytmetyce przedziałowej
	
	\begin{equation}
		p(x) - a_nx^n + a_{n-1}x^{n-1} + ... + a_1z + a_0
	\end{equation}
	
	\section{Opis metody}
	Do znalezienia pierwiastków zespolonych wielomianu (1) stosuje się metodę Bairstowa, w której unika się rachunków na liczbach zespolonych. Metoda ta opiera się na twierdzeniu, że pierwiastki wielomianu 
	
	\begin{equation}
		x^2 - rz - q
	\end{equation}

	gdzie r i q oznaczają współczynniki rzeczywiste, są pierwiastkami wielomianu (1), gdy wielomian (1) dzieli się bez reszty przez wielomian (2). W ogólności wykonując dzielenie wielomianu (1) przez wielomian (2) mamy
	
	\begin{equation}
		p(x) = p_1(x)(x^2 - rx - q) + Ax + B
	\end{equation}

	przy czym stopień wielomianu $p_1(x)$ jest nie większy niż $n - 2$, a wyrażenie $Ax + B$ jest resztą z dzielenia. Współczynniki A i B zależą od r i q:
	
	$$A = A(r,q), \quad B = B(r,q) $$
	
	Reszta z dzielenia będzie równa 0, jeżeli 
	
	$$ A(r,q) = 0, \quad B(r,q) = 0$$
	
		
		Powyższe równania są układem dwu równań nieliniowych o dwu niewiadomych $r$ i $q$. Do
		ich rozwiązania możemy zastosować proces iteracyjny Newtona:
		\[
		\begin{bmatrix}
			r_{j+1} \\
			q_{j+1}
		\end{bmatrix}
		=
		\begin{bmatrix}
			r_j \\
			q_j
		\end{bmatrix}
		-
		{\begin{bmatrix}
				\frac{\partial A}{\partial r} & \frac{\partial A}{\partial q} \\[1.5ex]
				\frac{\partial B}{\partial r} & \frac{\partial B}{\partial q}
		\end{bmatrix}}^{-1}_{\substack{r=r_j \\ q=q_j}}
		\begin{bmatrix}
			A(r_j, q_j) \\
			B(r_j, q_j)
		\end{bmatrix},
		\]
		
		który można zapisać w postaci układu równań liniowych:
		\begin{equation}
			{\begin{bmatrix}
					\frac{\partial A}{\partial r} & \frac{\partial A}{\partial q} \\[1.5ex]
					\frac{\partial B}{\partial r} & \frac{\partial B}{\partial q}
			\end{bmatrix}}_{\substack{r=r_j \\ q=q_j}}
			\begin{bmatrix}
				r_{j+1} \\
				q_{j+1}
			\end{bmatrix}
			=
			{\begin{bmatrix}
					\frac{\partial A}{\partial r} & \frac{\partial A}{\partial q} \\[1.5ex]
					\frac{\partial B}{\partial r} & \frac{\partial B}{\partial q}
			\end{bmatrix}}_{\substack{r=r_j \\ q=q_j}}
			\begin{bmatrix}
				r_j \\
				q_j
			\end{bmatrix}
			-
			\begin{bmatrix}
				A(r_j, q_j) \\
				B(r_j, q_j)
			\end{bmatrix}.
			\tag{4}
		\end{equation}
		
		Wartości $A(r, q)$ i $B(r, q)$ można otrzymać za pomocą następujących wzorów rekurencyjnych:
		\begin{gather*}
			b_0 = a_n, \quad b_1 = b_0 r + a_{n-1}, \\[1ex]
			b_i = b_{i-2} q + b_{i-1} r + a_{n-i}, \quad i = 2, 3, \dots, n-2, \\[1ex]
			A(r, q) = b_{n-3} q + b_{n-2} r + a_1, \quad B(r, q) = b_{n-2} q + a_0.
		\end{gather*}
		
		Pochodne cząstkowe występujące w równaniach (4) wyznacza się na podstawie wzorów
		\begin{gather*}
			\frac{\partial A(r, q)}{\partial r} = r\bar{A}(r, q) + \bar{B}(r, q), \quad \frac{\partial A(r, q)}{\partial q} = \bar{A}(r, q), \\[2ex]
			\frac{\partial B(r, q)}{\partial r} = q\bar{A}(r, q), \quad \frac{\partial B(r, q)}{\partial q} = \bar{B}(r, q),
		\end{gather*}
		
		gdzie wartości $\bar{A}(r, q)$ i $\bar{B}(r, q)$ mogą być obliczone na podstawie podobnych wzorów
		rekurencyjnych, a mianowicie
	
		\begin{gather*}
			c_0 = b_0, \quad c_1 = c_0r + b_1, \\
			c_i = c_{i-2}q + c_{i-1}r + b_i, \quad i = 2,3,\dots,n-4, \\
			\bar{A}(r,q) = c_{n-5}q + c_{n-4}r + b_{n-3}, \quad \bar{B}(r,q) = c_{n-4}q + b_{n-2}.
		\end{gather*}
	
	\section{Wywołanie funkcji}
	% Punkt 3: Instrukcja wywołania (jeśli funkcja) lub procedury

	\begin{lstlisting}[language=C++, frame=single, basicstyle=\ttfamily\small]
template<typename T>
vector<complex_interval<T>> bairstow_method(
	int degree,
	vector<T> coefficients,
	int max_iterations,
	long double tolerance,
	long double zero
);
\end{lstlisting}
	
	\section{Dane}
	% Punkt 4: Krótki opis parametrów wejściowych
	\begin{itemize}
		\item \texttt{int degree} -- stopień wielomianu.
		\item \texttt{vector<T> coefficients} -- zawierająca współczynniki wielomianu
		\item \texttt{int max\_iterations} -- liczba iteracji.
		\item \texttt{long double tolerance} -- parametr służący do przerwania iteracji
		\item \texttt{long double zero} -- wartość przyjmowana za zero maszynowe
	\end{itemize}
	% Uwaga: Jeśli indeksujesz od zera, zmień zakresy na 0...n-1 i opis na a_{i+1,j+1}
	
	\section{Wyniki}
	% Punkt 5: Opis parametrów wyjściowych
	\begin{itemize}
		\item \texttt{vector<complex\_interval<T>>} -- tablica dynamiczna zawierająca rozwiązanie, i-ty element tablicy zawiera pierwiastek zespolony wielomianu.
	\end{itemize}
	
	\section{Identyfikatory nielokalne}
	
	\begin{itemize}
		\item \texttt{complex\_interval<T>} -- nazwa typu strukturalnego reprezentująca dane zespolone o składowych typu \texttt{T}, zdefiniowanego następująco:
		\begin{lstlisting}[language=C++, basicstyle=\ttfamily\small]			
template<typename T>
struct complex_interval 
{
	T real;
	T imag;
};
		\end{lstlisting}
		
		\item \texttt{vector<T>} -- nazwa typu reprezentującego tablicę dynamiczną (kontener standardowy) o elementach typu \texttt{T}.
		
		\item \texttt{vector<complex\_interval<T>>} -- nazwa typu reprezentującego tablicę dynamiczną, której elementami są struktury przedziałów zespolonych.
	\end{itemize}
	
	\newpage
	
\section{Tekst funkcji}
% Punkt 9: Nalezy w tym miejscu zrobic copy & paste swojej funkcji 
\begin{lstlisting}[language=C++, frame=single, basicstyle=\ttfamily\small, breaklines=true, columns=fullflexible, keepspaces=true]
template<typename T>
vector<complex_interval<T>> bairstow_method(
	int degree, 
	vector<T> coefficients, 
	int max_iterations, 
	long double tolerance, 
	long double zero)
{
	vector<complex_interval<T>> output;
	
	if (degree == 0) return output;
		
	if (degree == 1) {
		return first_degree_roots<T>(coefficients, zero);
	}
		
	if (degree == 2) {
		return second_degree_roots<T>(coefficients, zero);
	}
		
	newton_output result = newton<T>(coefficients,max_iterations,tolerance);
		
	T a_coef;
	if constexpr (is_same_v<T, long double>) {
		a_coef = 1.0;
	} else if constexpr (is_same_v<T, Interval<long double>>) {
		a_coef = IntRead<long double>("1.0");
	}
		
	// Obliczanie pierwiastkow z wyznaczonego czynnika kwadratowego
	vector<complex_interval<T>> temp = second_degree_roots<T>(
	{-1 * result.q, -1 * result.r, a_coef}, zero);
	output.insert(output.end(), temp.begin(), temp.end());
		
	// Rekurencyjne wywolanie dla wielomianu o stopniu n-2
	vector<complex_interval<T>> rest = bairstow_method<T>(
	static_cast<int>(result.coefficients.size() - 1), 
	result.coefficients, 
	max_iterations, 
	tolerance, 
	zero);
	output.insert(output.end(), rest.begin(), rest.end());
		
	return output;
	
	}
\end{lstlisting}
	


\section{Przykłady}

\subsection*{a) Dane rzeczywiste dla arytmetyki zwykłej}
\begin{itemize}
	\item \textbf{Wielomian:} $P(x) = x^4 + 5x^3 + 10x^2 + 10x + 4$ \\
	\textbf{Dane:} \texttt{1,4 5,3 10,2 10,1 4,0} \\
	\textbf{Wyniki:}
	\begin{itemize}
		\item $x_1 = \texttt{-1.0000000000000000E0 + 0.0000000000000000E-1i}$
		\item $x_2 = \texttt{-2.0000000000000001E0 + 0.0000000000000000E-1i}$
		\item $x_3 = \texttt{-9.9999999999999992E-1 + 1.0000000000000000E0i}$
		\item $x_4 = \texttt{-9.9999999999999992E-1 - 1.0000000000000000E0i}$
	\end{itemize}
	
	\item \textbf{Wielomian:} $P(x) = x^2 - 5x + 6$ \\
	\textbf{Dane:} \texttt{1,2 -5,1 6,0} \\
	\textbf{Wyniki:}
	\begin{itemize}
		\item $x_1 = \texttt{3.0000000000000000E0 + 0.0000000000000000E-1i}$
		\item $x_2 = \texttt{2.0000000000000000E0 + 0.0000000000000000E-1i}$
	\end{itemize}
\end{itemize}

\subsection*{b) Dane rzeczywiste dla arytmetyki przedziałowej}
\begin{itemize}
	\item \textbf{Wielomian:} $P(x) = x^5 - 15x^4 + 85x^3 - 225x^2 + 274x - 120$ \\
	\textbf{Dane:} \texttt{1,5 -15,4 85,3 -225,2 274,1 -120,0} \\
	\textbf{Wyniki:}
	\begin{itemize}
		\item $x_1 = \texttt{[4.0000000000000000E0, 4.0000000000000001E0]}\quad
		$\texttt{w: 0.00E-1}$ \ +\\ \, \texttt{[0.0000000000000000E-1, 0.0000000000000000E-1]i}$ \quad
		$\texttt{w: 0.00E-1}$
		
		\item $x_2 = \texttt{[9.9999999999999999E-1, 1.0000000000000000E0]} \quad
		$\texttt{w: 0.00E-1}$ \ +\\ \, \texttt{[0.0000000000000000E-1, 0.0000000000000000E-1]i}$ \quad
		$\texttt{w: 0.00E-1}$
		
		\item $x_3 = \texttt{[2.9999999999999999E0, 3.0000000000000001E0]} \quad
		$\texttt{w: 0.00E-1}$ \ +\\ \, \texttt{[0.0000000000000000E-1, 0.0000000000000000E-1]i}$ \quad
		$\texttt{w: 0.00E-1}$
		
		\item $x_4 = \texttt{[1.9999999999999999E0, 2.0000000000000000E0]} \quad
		$\texttt{w: 0.00E-1}$ \ +\\ \, \texttt{[0.0000000000000000E-1, 0.0000000000000000E-1]i}$ \quad
		$\texttt{w: 0.00E-1}$
		
		\item $x_5 = \texttt{[5.0000000000000003E0, 5.0000000000000004E0]} \quad
		$\texttt{w: 0.00E-1}$ \ +\\ \, \texttt{[0.0000000000000000E-1, 0.0000000000000000E-1]i}$ \quad
		$\texttt{w: 0.00E-1}$
	\end{itemize}

	\newpage
	
	\item \textbf{Wielomian:} $P(x) = x^6 - 7x^4 + 14x^2 - 8$ \\
	\textbf{Dane:} \texttt{1,6 0,5 -7,4 0,3 14,2 0,1 -8,0} \\
	\textbf{Wyniki:}
	
	\begin{itemize}
		\item $x_1 = \texttt{[1.0000000000000000E0, 1.0000000000000000E0]} \quad \texttt{w: 0.00E-1} \ + \, \texttt{[0.0000000000000000E-1, 0.0000000000000000E-1]i} \quad \texttt{w: 0.00E-1}$
		
		\item $x_2 = \texttt{[-1.0000000000000000E0, -1.0000000000000000E0]} \quad \texttt{w: 0.00E-1} \ + \, \texttt{[0.0000000000000000E-1, 0.0000000000000000E-1]i} \quad \texttt{w: 0.00E-1}$
		
		\item $x_3 = \texttt{[1.4142135623730950E0, 1.4142135623730951E0]} \quad \texttt{w: 0.00E-1} \ + \, \texttt{[0.0000000000000000E-1, 0.0000000000000000E-1]i} \quad \texttt{w: 0.00E-1}$
		
		\item $x_4 = \texttt{[-1.4142135623730951E0, -1.4142135623730950E0]} \quad \texttt{w: 0.00E-1} \ + \, \texttt{[0.0000000000000000E-1, 0.0000000000000000E-1]i} \quad \texttt{w: 0.00E-1}$
		
		\item $x_5 = \texttt{[1.9999999999999999E0, 2.0000000000000000E0]} \quad \texttt{w: 0.00E-1} \ + \, \texttt{[0.0000000000000000E-1, 0.0000000000000000E-1]i} \quad \texttt{w: 0.00E-1}$
		
		\item $x_6 = \texttt{[-2.0000000000000000E0, -1.9999999999999999E0]} \quad \texttt{w: 0.00E-1} \ + \, \texttt{[0.0000000000000000E-1, 0.0000000000000000E-1]i} \quad \texttt{w: 0.00E-1}$
	\end{itemize}
	
\end{itemize}

	
	
	\subsection*{c) Dane przedziałowe dla arytmetyki przedziałowej}
	\begin{itemize}
		
	
	\item \textbf{Wielomian:} $P(x) = [0.99,1.01]x^2 - [4.99,5.01]x + [5.99,6.01]$ \\
	\textbf{Dane:} \texttt{[0.99,1.01],2  [-5.01,-4.99],1 [5.99,6.01],0} \\
	\textbf{Wyniki:}
	\begin{itemize}
		\item $x_1 = \texttt{[2.9177829616973363E0, 3.1235385439530877E0]} \quad \texttt{w: 2.10E-1} \ +\\ \, \texttt{[0.0000000000000000E-1, 0.0000000000000000E-1]i} \quad \texttt{w: 0.00E-1}$
		
		\item $x_2 = \texttt{[1.9269665065519628E0, 2.1327220888077142E0]} \quad \texttt{w: 2.10E-1} \ +\\ \, \texttt{[0.0000000000000000E-1, 0.0000000000000000E-1]i} \quad \texttt{w: 0.00E-1}$
	\end{itemize}
\end{itemize}

\subsection*{d) Przypadek nieobsługiwany}
*Algorytm nie radzi sobie z wielomianami z pierwiastkami wielokrotnymi
\begin{itemize}
	\item \textbf{Wielomian:} $P(x) = x^6 - 6x^5 + 15x^4 - 20x^3 + 15x^2 - 6x + 1$ \\
	\textbf{Dane:} \texttt{1,6 -6,5 15,4 -20,3 15,2 -6,1 1,0} \\
	\textbf{Wyniki:}
	\begin{itemize}
		\item $x_1 = \texttt{[1.0004780799214135E0, 1.0004780799214137E0]} \quad \texttt{w: 0.00E-1} \ + \, \texttt{[0.0000000000000000E-1, 0.0000000000000000E-1]i} \quad \texttt{w: 0.00E-1}$
		
		\item $x_2 = \texttt{[9.9952233407555942E-1, 9.9952233407555954E-1]} \quad \texttt{w: 0.00E-1} \ + \, \texttt{[0.0000000000000000E-1, 0.0000000000000000E-1]i} \quad \texttt{w: 0.00E-1}$
		
		\item $x_3 = \texttt{[9.9956389763679445E-1, 9.9956389763679446E-1]} \quad \texttt{w: 0.00E-1} \ + \, \texttt{[7.5470439697661299E-4, 7.5470439697664891E-4]i} \quad \texttt{w: 0.00E-1}$
		
		\item $x_4 = \texttt{[9.9956389763679445E-1, 9.9956389763679446E-1]} \quad \texttt{w: 0.00E-1} \ + \, \texttt{[-7.5470439697664891E-4, -7.5470439697661299E-4]i} \quad \texttt{w: 0.00E-1}$
		
		\item $x_5 = \texttt{[1.0004357292609590E0, 1.0004357292609591E0]} \quad \texttt{w: 0.00E-1} \ + \, \texttt{[7.5535062820352275E-4, 7.5535062820502989E-4]i} \quad \texttt{w: 0.00E-1}$
		
		\item $x_6 = \texttt{[1.0004357292609590E0, 1.0004357292609591E0]} \quad \texttt{w: 0.00E-1} \ + \, \texttt{[-7.5535062820502989E-4, -7.5535062820352275E-4]i} \quad \texttt{w: 0.00E-1}$
	\end{itemize}

\end{itemize}
	
	
\end{document}